% Clase 14 --- No hay monopolos magnéticos (§2.3).
% "Agarro una buena sierra y quiero separar el polo norte del polo sur": el
% experimento mental que el docente dibujó. El resultado no es un polo suelto
% sino dos imanes completos, y por más que se repita el corte sigue pasando lo
% mismo.
\begin{center}
\begin{tikzpicture}[scale=1]

  % =============== el imán entero ===============
  \begin{scope}
    \fill[figblue!18] (-1.6,-0.45) rectangle (0,0.45);
    \fill[figred!18] (0,-0.45) rectangle (1.6,0.45);
    \draw[figgray, line width=1pt] (-1.6,-0.45) rectangle (1.6,0.45);
    \draw[figgray, line width=0.7pt] (0,-0.45) -- (0,0.45);
    \node[figlbl, figblue] at (-0.8,0) {S};
    \node[figlbl, figred] at (0.8,0) {N};

    % la sierra
    \draw[figgray, line width=0.8pt] (0,0.75) -- (0,1.25);
    \node[figlblaux, anchor=south] at (0,1.28) {cortar acá};

    \node[figlbl, anchor=north] at (0,-0.95) {un imán};
  \end{scope}

  % ---- flecha ----
  \draw[figvecaux, line width=1pt] (2.3,0) -- (3.5,0);
  \node[figlblaux, anchor=south] at (2.9,0.1) {cortar};

  % =============== los dos pedazos ===============
  \begin{scope}[xshift=6.1cm]
    % pedazo izquierdo
    \fill[figblue!18] (-1.9,-0.45) rectangle (-1.1,0.45);
    \fill[figred!18] (-1.1,-0.45) rectangle (-0.3,0.45);
    \draw[figgray, line width=1pt] (-1.9,-0.45) rectangle (-0.3,0.45);
    \draw[figgray, line width=0.7pt] (-1.1,-0.45) -- (-1.1,0.45);
    \node[figlbl, figblue, scale=0.9] at (-1.5,0) {S};
    \node[figlbl, figred, scale=0.9] at (-0.7,0) {N};

    % pedazo derecho
    \fill[figblue!18] (0.3,-0.45) rectangle (1.1,0.45);
    \fill[figred!18] (1.1,-0.45) rectangle (1.9,0.45);
    \draw[figgray, line width=1pt] (0.3,-0.45) rectangle (1.9,0.45);
    \draw[figgray, line width=0.7pt] (1.1,-0.45) -- (1.1,0.45);
    \node[figlbl, figblue, scale=0.9] at (0.7,0) {S};
    \node[figlbl, figred, scale=0.9] at (1.5,0) {N};

    \node[figlbl, figred, anchor=north, align=center] at (0,-0.95)
      {dos imanes completos,\\[-0.2em] no dos polos sueltos};
  \end{scope}

  % ---- y otra vez ----
  \draw[figvecaux, line width=1pt] (8.4,0) -- (9.4,0);
  \node[figlblaux, anchor=south, align=center] at (8.9,0.1) {y otra vez};
  \node[figlbl, anchor=west, align=left] at (9.6,0)
    {\dots\ siempre igual};

\end{tikzpicture}\\[0.5em]
{\small\itshape El experimento se puede repetir todas las veces que se quiera y
el resultado no cambia: nunca queda un polo aislado. Esto no es un detalle
técnico sino una diferencia estructural con el caso eléctrico, donde sí se puede
tener una carga positiva sola, y es la razón de fondo por la que el magnetismo
no se puede describir copiando la electrostática: si no hay monopolos, el papel
de \enquote{fuente} del campo tiene que jugarlo otra cosa, y esa otra cosa son
las \textbf{cargas en movimiento}. Como curiosidad, algunas teorías predicen
monopolos, y bastaría con que existieran unos pocos para que quedara explicado
por qué la carga eléctrica está cuantizada ---pero pese a muchísima búsqueda
experimental no se detectó ninguno hasta hoy---.}
\end{center}
